\documentclass{jsarticle}
\usepackage{amsmath}
\usepackage{amssymb}
\begin{document}
\title{微分幾何の量子化とゼータ関数、\\
大域的微分方程式についてのレポート}
\author{Masaaki Yamaguchi, with my son}
\date{}
\maketitle

  $$ \text{ヒッグス場 + オイラーの定数　= ゼータ関数} $$
  $$ m(x) + C = 2 {d \over df}\int\int {1 \over (y \log y)^{1 \over 2}}dy_m$$
  $$ \text{オイラーの定数} = {d \over d \gamma}\Gamma^{-1} - \gamma^{(\gamma)
  ^{'}} = {\left(\Gamma^{(\gamma)^{'}}\right)}^{-1} - \gamma^{\gamma^{'}} $$
  $$ \text{重力 + 反重力 = ゼータ関数} $$
  ガンマ関数の大域的微分方程式が、大域的微分方程式の部分積分として、
  ヘルマンダー作用素と同型になり、微分の変分法が対数関数を初期関数
  として、物体を形作っているフェルミオンを商代数で求めると、
  ヒッグス場の方程式になる。
  $$ \Gamma dx_m + \int \Gamma dx_m = \Gamma^{(\gamma)^{'}} + \Gamma^{(
  \gamma)} $$　
  $$ {d \over d \gamma}\Gamma + \int \Gamma dx_m $$
  $$ = \int f(g(x))g^{'}(x)dx $$
  $$ = \left(\int f(g(x))dx\right)^{'} $$
  $$ \nabla_i \nabla_j \int \nabla f(x) d \eta $$
  大域的微分方程式の部分積分をフロベニウスの定理で解を出すと、
  前後の大小関係が、Jones多項式から、rico level理論となり、
  ゼータ関数の対称率が、CP対称性の破れとして、宇宙のゼータ関数
  だけが出力される。これが、オイラーの定数の関係として、
  ゼータ関数とオイラーの定数の対称率の変分率としての、
  Jones多項式として、微分幾何の量子化が、ゼータ関数と証明されている。
  $$ |f \circ g| \le |f + g| $$
  $$ {f \over \log x} = m(x) $$
  $$ {\gamma \over \log x} = {d \over d \gamma}\Gamma $$
  $$ \gamma = \Gamma^{(\gamma)^{'}}\log x$$
  $$ \Gamma^{(\gamma)^{'}} = \gamma (\log x)^{-1} $$
  $$ = r dx_m $$
  $$ \Gamma^{(\gamma)^{'}} = \left(\int e^{-x}x^{1 - t}dx\right)^{
	  \left(\int e^{-x}x^{1 - t} \log x dx\right)^{'}} $$
  $$ = e^{f} $$
  $$ = r dx_m = e^{f} $$
 
  $$ \int \Gamma(\gamma)^{'}dx_m $$
  $$ \nabla_i \nabla_j \int \nabla \Gamma(\gamma)dx $$
  $$ = \int \Gamma dx_m \cdot {d \over d \gamma}\Gamma dx_m $$
  $$ = \int \Gamma \cdot {d \over d \gamma}\Gamma dx_m $$
  $$ = \left(\int \Gamma dx_m \circ {d \over d\gamma}\Gamma \le
  \int \Gamma dx_m + \Gamma dx_m \right)^{'} 
  = e^{-f} \cdot e^{f} \le e^{-f} + e^{f} $$
  $$ \left(1 \le e^{-f} + e^{f}\right)^{'} $$
  $$ 0 = -(e^{-f} - e^{f}) $$
  $$ \Gamma^{(\gamma)^{'}} + \Gamma^{(\gamma)} $$
  $$ (= e^f + e^{-f})^{'} $$
  $$ = - e^{-f} + e^{f} $$
  $$ = - (e^{-f} - e^f) $$
  $$ (R_{ij})^{'} = - (R_{ij}) $$
  $$ \int C dx_m = \int\left(\int{1 \over x^s}dx - \log x \right)dx $$
  $$ = \int \int{1 \over x^s}\mathrm{dvol} - \int \log x \mathrm{dvol} $$
  $$ \Gamma dx_m = \gamma dx_m $$
  $$ = {d \over d \gamma}\Gamma^{^1} - (\gamma)^{\gamma^{'}} $$
  $$ = e^{-f} - e^{f} $$
  微分幾何の量子化が、ガンマ関数となり、ベータ関数に行き着き、
  オイラーの定数へと導かれて、ヒッグス場と対になり、ゼータ関数へと
  証明されて、ポワンカレ予想とリーマン予想が解決されている。
  ロシアとアメリカ、日本の数学者たちが解決している。
  私の子どもが、微分幾何の量子化がゼータ関数へと解決している。
  全部の理論がヒルベルト空間を多様体として、一般相対性理論
  と量子力学を、微分幾何の量子化が、D-braneと解として、
  ガウスの曲面論を、重力と原子にして、統一場理論として形成している。
\end{document}
